%!TEX root = ../main.tex
\chapter{Background}
\label{ch:background}

This chapter sets out the way an autonomous robotic system is conventionally
built. It is not a survey: it is the minimum vocabulary needed to read the
chapters that follow, and the statement of the problem they address.

\section{The Node Model of ROS~2}
\label{sec:bg-ros2}

A robot is not a program. It is a set of concurrent processes that read sensors
at different rates, estimate quantities, decide and command actuators, and that
must keep working when one of them stops. Robotic \emph{middleware} is the
infrastructure that makes it possible to compose a system of this kind without
its components knowing one another.

The model adopted by ROS~2 \cite{macenski2022ros2} organises the system as a
graph of \emph{nodes}, independent processes that communicate by publishing and
subscribing typed messages on named channels. A node that publishes does not
know who reads it, nor whether anyone reads it; a node that subscribes does not
know who produces the datum. The coupling between components is therefore
limited to the type of the message and the name of the channel, and from this
property follow the three that matter here: a component can be replaced without
modifying the others; the system can be observed from the outside by listening
to the channels, which makes diagnosis a measurement rather than a conjecture;
and one and the same node can operate indifferently on data produced by a
simulator or by a real device, provided they arrive in the same form.

To this is added a mechanism specific to geometry. The measurements of the
sensors are each expressed in their own reference frame, and the middleware
maintains a tree of the transforms between such frames, updated over time, which
any node may query in order to carry a quantity from one frame to another. It is
the infrastructure on which any form of fusion between different sensors rests.

\section{The Classical Autonomy Stack}
\label{sec:bg-stack}

On that infrastructure, established practice builds a chain of stages, each of
which consumes the product of the preceding one.

\paragraph{Perception.}
The raw sensor data are validated, filtered and converted into a common form. It
is the stage at which what constitutes a valid measurement is decided, and in
which the characteristics of the individual device --- range, field of view,
noise, rate --- are taken into account.

\paragraph{State estimation.}
The measurements that observe motion are combined into a single estimate of the
pose of the vehicle. The standard instrument is a recursive Bayesian filter, in
its extended or unscented variant, which maintains the state and corrects it at
every incoming measurement, weighing it by its own uncertainty
\cite{moore2014localization}. Since every source observes only a part of the
motion and each drifts in its own way, the quality of the result depends on how
faithfully each declares its own uncertainty.

\paragraph{Map construction.}
The pose estimate on its own drifts: the error accumulates and nothing in the
estimation chain can notice it. Simultaneous localisation and mapping corrects
that drift by recognising places already visited and redistributing the
correction over the path \cite{cadena2016slam}. Alongside the global map so
obtained, navigation requires a local representation of free and occupied space,
typically an occupancy grid, planar or volumetric \cite{hornung2013octomap}.

\paragraph{Planning and control.}
On the representation so constructed, a planner computes a path towards the goal
and a controller translates it into velocity commands that respect the kinematic
limits of the platform. Mature navigation stacks \cite{macenski2020nav2}
organise this part as a state machine alternating global planning, local
tracking and recovery behaviours when the path is not traversable.

\medskip

The scheme has a notable property: every stage is inspectable and every decision
is traceable to a rule someone wrote. It is the reason why it remains the
industrial standard, and why in degraded environments it still represents the
highest engineering level reached, as the heterogeneous robot teams of the DARPA
Subterranean Challenge have shown \cite{tranzatto2022cerberus}. It is also the
reason why this thesis does not replace it in its entirety.

\section{The Cost of Heterogeneity}
\label{sec:bg-portability}

The scheme has a cost, however, which appears when the same system must work on
different platforms. Each stage is parameterised, and its parameters are not
universal properties: they are tunings. The threshold beyond which a range
measurement is considered noise depends on the sensor; the parameters of the
estimation filter depend on which sources are present and on how reliable each
of them is on that platform; the inflation radius of the obstacles depends on
the extent of the vehicle; the gains of the controller depend on the dynamics of
the body.

It follows that what survives a change of robot is the structure of the chain,
not the working system. Porting a proven stack to a new platform means going
through every stage and retuning it, and the cost recurs at every variation of
the sensor suite. The problem is the more relevant the more heterogeneous the
fleet: the case treated in this thesis comprises ground and aerial platforms
whose masses differ by two orders of magnitude --- a rover of \SI{200}{\kilo\gram}
and a multirotor of under \SI{2}{\kilo\gram} --- one of which has a degree of
freedom the other lacks, and with different sets of sensors.

A second cost concerns the decision stage. Classical planning requires the
environment to be represented in a form on which a rule can operate, and
requires the rule to anticipate the cases it will encounter. Where the
environment is cluttered, dynamic or never mapped, the set of anticipated cases
proves incomplete, and the resulting behaviour is conservative at best and
deadlocked at worst.

It is on these two costs that the present work acts, and it acts on them
asymmetrically: the first four stages are preserved, but generated from a
declaration of the sensors present instead of being tuned for one platform
(Chapter~\ref{ch:pipeline}), while the decision stage is replaced by a learnt
policy that reads the local representation of space, the pose and the goal, and
produces the velocity command directly (Chapters~\ref{ch:lnn}
and~\ref{ch:training}).
